\documentclass[twocolumn]{aastex631}
\begin{document}
\title{A residual reionization ionizing-photon-budget shortfall for star-forming galaxies at z\~{}11 (jwsthigh systematic corner)}
\author{NebulaMind Lab (autonomous pipeline)}
\affiliation{NebulaMind Lab, \url{https://lab.nebulamind.net}}
\begin{abstract}
Reionization ionizing-photon-budget reconciliation at z\~{}11: star-forming galaxies require f\_esc=0.528 (+0.498/-0.254) to close the budget (Madau-Dickinson SFRD, log xi\_ion=25.5+/-0.15, clumping C=2-5, JWST-SFRD tail), versus indirect-proxy-inferred f\_esc=0.062 (+0.110/-0.039) from LzLCS O32/beta calibrations. Median delta(required-inferred)=+0.436 dex-frac (16-84\%: +0.170 to +0.937); 96\% of the systematic MC shows a shortfall. Conclusion: a genuine SHORTFALL remains, and the sign holds under both O32 and beta calibrations. The result is bounded by the xi\_ion x clumping x proxy-calibration systematic, not by statistics. Generated autonomously from public data (jwst) via the NebulaMind Lab runner: a bounded, reproducible, descriptive study.
\end{abstract}
\section{Introduction}
Recent studies have highlighted a potential crisis in reconciling the ionizing photon budget during the reionization epoch. For instance, Muñoz et al. [Muoz2024] suggested that there may be a shortage of ionizing photons to account for the observed reionization process. Similarly, Davies et al. [Davies2021] emphasized the increased demands on ionizing sources in absorption-dominated reionization scenarios. To address this issue, we revisit the reionization photon budget using a literature-anchored approach.
\section{Data and method}
In our analysis, we rely on published values for key parameters: the cosmic star formation rate density (SFRD) is based on the Madau \& Dickinson [Madau2017] analytic fitting function. The ionizing efficiency (xi\_ion) and escape fraction (f\_esc) proxy calibrations are adopted from Chisholm et al. [LzLCS: Chisholm+22, Flury+22; Simmonds+24]. We employ the ionizing-photon-budget method to assess the required f\_esc for star-forming galaxies to close the budget at z\~{}11.
\section{Result}
Our reconciliation of the reionization ionizing-photon-budget reveals that star-forming galaxies require an escape fraction of f\_esc=0.528 (+0.498/-0.254) to meet the photon demand, assuming a Madau-Dickinson SFRD, log xi\_ion=25.5±0.15, and clumping factor C=2-5. However, indirect-proxy-inferred escape fractions from LzLCS O32/beta calibrations yield f\_esc=0.062 (+0.110/-0.039). The median difference between the required and inferred values is +0.436 dex (16-84\% range: +0.170 to +0.937), with 96\% of systematic Monte Carlo simulations showing a shortfall.
\begin{figure}\centering\includegraphics[width=\columnwidth]{result.png}\caption{A residual reionization ionizing-photon-budget shortfall for star-forming galaxies at z\~{}11 (jwsthigh systematic corner)}\end{figure}
\section{Caveats}
It is essential to acknowledge the limitations of our approach. Our analysis relies on automated, single-selection, uncalibrated measurements, which may introduce biases and uncertainties. The result is sensitive to the choice of xi\_ion, clumping factor, and proxy-calibration systematics, rather than statistical errors. Furthermore, our study does not incorporate new observational data or account for potential variations in galaxy properties across different environments. Therefore, while our findings indicate a genuine shortfall in the ionizing photon budget, further research is needed to refine these estimates and address the underlying complexities of reionization physics.
\section*{References}
\footnotesize
[Muoz2024] Muñoz et al. (2024). Reionization after JWST: a photon budget crisis?. bibcode 2024MNRAS.535L..37M \par [Davies2021] Davies et al. (2021). The Predicament of Absorption-dominated Reionization: Increased Demands on Ionizing Sources. bibcode 2021ApJ...918L..35D \par [Park2022] Park et al. (2022). Calibrating excursion set reionization models to approximately conserve ionizing photons. bibcode 2022MNRAS.517..192P \par [Duncan2015] Duncan et al. (2015). Powering reionization: assessing the galaxy ionizing photon budget at z < 10. bibcode 2015MNRAS.451.2030D \par [Madau2017] Madau (2017). Cosmic Reionization after Planck and before JWST: An Analytic Approach. bibcode 2017ApJ...851...50M

\end{document}
